\documentclass[11pt]{article}
\usepackage[T1]{fontenc}
\usepackage[utf8]{inputenc}
\usepackage{lmodern}
\input{glyphtounicode}
\pdfgentounicode=1
\usepackage{amsmath,amssymb,amsthm,geometry,hyperref}
\geometry{margin=1in}
\emergencystretch=1em
\hypersetup{colorlinks=true,linkcolor=blue,citecolor=blue,urlcolor=blue,pdftitle={A counterexample to Stanley's conjecture on Smith forms of Jucys--Murphy elements},pdfauthor={P. Graham},pdfsubject={Smith normal forms of Jucys--Murphy elements on integral Specht modules},pdfkeywords={Smith normal form, Jucys--Murphy element, Specht module, symmetric group, Lean}}
\newtheorem{theorem}{Theorem}
\newtheorem{lemma}[theorem]{Lemma}
\newtheorem{corollary}[theorem]{Corollary}
\newtheorem{proposition}[theorem]{Proposition}
\theoremstyle{remark}
\newtheorem{remark}[theorem]{Remark}
\newcommand{\Z}{\mathbb{Z}}
\newcommand{\Q}{\mathbb{Q}}
\newcommand{\F}{\mathbb{F}}
\newcommand{\sn}{\mathfrak{S}_n}
\newcommand{\snf}{\stackrel{\mathrm{snf}}{\rightarrow}}
\newcommand{\file}[1]{\texttt{\detokenize{#1}}}
\newcommand{\tb}[3]{\begin{array}{@{}l@{}}#1\\[-2pt]#2\\[-2pt]#3\end{array}}
\title{A counterexample to Stanley's conjecture on Smith forms of Jucys--Murphy elements}
\author{P. Graham\\ \small\url{https://orcid.org/0009-0007-8493-1505}}
\date{October 21, 2026}
\begin{document}
\maketitle
\begin{abstract}
For a partition $\lambda$ of $n$ with $f^\lambda$ standard Young tableaux,
let $S^\lambda_\Z$ be the integral Specht module and
$X_k=(1,k)+\cdots+(k-1,k)$ the $k$-th Jucys--Murphy element. Stanley
conjectured in 2016 that the invariant factors of $X_k$ on $S^\lambda_\Z$,
taken nonnegative, are $\alpha_{f^\lambda-r+1}=|\prod(j-i)|$ for
$1\le r\le f^\lambda$, the product (empty product $1$) over the cells
$(i,j)$ that $k$ occupies in at least $r$ standard tableaux. We show that it
is false. In $\Z\sn$, $X_3\bigl(X_3-(1,2)\bigr)=2$. Hence, for $n\ge3$, on
every full-rank $\Z$-lattice in $S^\lambda_\Q$ stable under the permutations
of $\{1,2,3\}$, including $S^\lambda_\Z$, $X_3$ has Smith form
$(1^{f^\lambda-m},2^m)$, where $m$ counts the tableaux with $3$ in cell
$(1,3)$ or $(3,1)$. So the conjecture holds at $k=3$ exactly when
$\lambda_1\le2$ or $\lambda$ has at most two rows; the smallest failure,
$\lambda=(3,1,1)$, gives $(1,1,1,1,2,2)$ against the predicted
$(1,1,1,1,1,4)$, and no case with $n\le4$ fails. Lean checks the identity
and the $6\times6$ matrix computation for $(3,1,1)$, independent programs
check that this matrix represents $X_3$, and the general argument is a
written proof.
\end{abstract}

\noindent The proofs, the Lean formalization, the certificate checkers and the text were produced by AI models (from OpenAI and Anthropic), working under the author's direction.

\medskip\noindent Status: Computer-assisted; not yet reviewed by experts.

\section{Introduction}

In his survey \emph{Smith normal form in combinatorics}~\cite{stanley},
Richard P.~Stanley proposed studying the Smith normal form (SNF) of
natural elements of the group ring $\Z\sn$. His example is the
\href{https://en.wikipedia.org/wiki/Jucys\%E2\%80\%93Murphy_element}{Jucys--Murphy
element} $X_k=(1,k)+(2,k)+\cdots+(k-1,k)$ \cite{jucys,murphy}, acting on the
Specht module $S^\lambda$, where its eigenvalues are the contents of the
cells occupied by $k$ in the standard Young tableaux of shape $\lambda$. The
survey reports that Darij Grinberg computed the SNF of the matrices
$W_{\lambda,k}$ for $\lambda\vdash n\le 7$, and on the basis of these data
Stanley formulated Conjecture~7.4: the Smith form is given by a product
formula in the contents, recalled exactly in Section~\ref{sec:statement}.
Stanley illustrated it with $\lambda=(5,1)$, $k=5$, where the formula gives
$(1,1,3,3,12)$.

On December 25, 2024, Alexander Chervov asked on MathOverflow about the
status of these conjectures~\cite{mo}; the question had no answer as of
October 21, 2026. In a comment on it, Grinberg noted that the survey defines
the matrices $W_{\lambda,k}$ ``somewhat clumsily'', through a basis of
$\Z\sn$, and clarified the intended lattice: the integral Specht module with
its standard polytabloid basis. We use exactly this reading.

Our main observation is an identity in $\Z\mathfrak{S}_3\subseteq\Z\sn$
(Lemma~\ref{lem:identity}):
\[
X_3\bigl(X_3-(1,2)\bigr)=\bigl(X_3-(1,2)\bigr)X_3=2 .
\]
It forces every invariant factor of $X_3$ to divide $2$, on every lattice
stable under the permutations of $\{1,2,3\}$.

\begin{theorem}\label{thm:k3}
Let $\lambda\vdash n\ge3$, and let $m$ be the number of standard Young
tableaux of shape $\lambda$ in which $3$ occupies the cell $(1,3)$ or $(3,1)$.
On every full-rank $\Z$-lattice in $S^\lambda_\Q$ that is stable under the
permutations of $\{1,2,3\}$, in particular on $S^\lambda_\Z$, the Smith
normal form of $X_3$ is $(1^{f^\lambda-m},2^m)$.
\end{theorem}

\begin{corollary}\label{cor:k3}
Conjecture~7.4 holds at $k=3$ if and only if $\lambda_1\le2$ or $\lambda$
has at most two rows. When $\lambda_1\ge3$ and $\lambda$ has at least three
rows, it fails on every lattice as in Theorem~\ref{thm:k3}: the predicted
largest invariant factor is $4$, while every actual one divides $2$.
\end{corollary}

The smallest such shape is $\lambda=(3,1,1)$. There the Smith form of $X_3$
on $S^{(3,1,1)}_\Z$ is $(1,1,1,1,2,2)$, while Conjecture~7.4 predicts
$(1,1,1,1,1,4)$; Section~\ref{sec:witness} gives the $6\times6$ matrix
explicitly. No case with $n\le4$ fails, and at $n=5$ the only failures are
$\lambda=(3,1,1)$ with $k=3$ and with $k=5$ (Proposition~\ref{prop:small});
these are the smallest counterexamples in this sense. A census of all
$\lambda\vdash n\le 8$ and all $k$ finds 94 disagreements among 416 cases
(Section~\ref{sec:census}).

\paragraph{What is checked, and how.} Lean proves the identity of
Lemma~\ref{lem:identity} in the group ring and, for the explicit $6\times6$
integer matrix $A$ of Section~\ref{sec:witness}, its Smith form, the relation
$AB=2I$ with an explicit integer matrix $B$, and that the predicted form is
not reachable from $A$. That $A$ is the matrix of $X_3$ on $S^{(3,1,1)}_\Z$
in the standard polytabloid basis is established by two independent
certificate checkers, each working from the definitions of tableaux,
tabloids and polytabloids rather than trusting the encoded matrix: one
recomputes $A$, the other verifies the defining equations
$X_3e_{T_j}=\sum_iA_{ij}e_{T_i}$ in the tabloid module, which determine $A$.
Theorem~\ref{thm:k3}, Corollary~\ref{cor:k3} and the hand argument for
$n\le4$ are written proofs and are not machine-checked; a third checker
recomputes every case with $n\le5$. Checker~A also verifies, for this
$\lambda$, that the standard polytabloids form a $\Z$-basis of the span of all
polytabloids, so the witness does not rest on the standard basis theorem; the
two checkers enumerate the standard tableaux by different methods. The census
for $n\le8$ (Sections~\ref{sec:readings}--\ref{sec:census}) is supporting
data, not part of any proof: it is computed by a construction that uses no
standard basis and agrees with a second, standard-basis computation for all
$n\le6$; its programs and output are included. Section~\ref{sec:verify}
gives details.

\paragraph{Credits.} The conjecture, its formulation and its worked example
are Stanley's. The computations on which it rests are reported in the survey
as Grinberg's, and the clarification of the intended lattice is his.
Chervov's question drew attention to the status of the conjecture. We know
of no earlier report that the conjecture fails. We do not know why the
computations reported in the survey supported the conjecture, and we
attribute no error to them (Section~\ref{sec:open}).

\section{The statement and the reading used}\label{sec:statement}

The following is quoted from the final part of Section~7 of~\cite{stanley}
(arXiv version~2, page~15; the text of version~1 is the same). Reference
numbers inside quotations are those of the quoted source.

\begin{quote}\sloppy
One intriguing example is the \emph{Jucys-Murphy element} (though it first
appears in the work of Alfred Young [50, \S19]) $X_k\in\Q\sn$,
$1\leq k\leq n$. It is defined by $X_1=0$ and
\[ X_k = (1,k)+(2,k)+\cdots+(k-1,k),\ \ 2\leq k\leq n, \]
where $(i,k)$ denotes the transposition interchanging $i$ and $k$. Just as
for $\Gamma_n(q)$, we can choose an integral basis for $\Q\sn$ (that is, a
$\Z$-basis for $\Z\sn$) so that the action of $X_k$ on $\Q\sn$ with respect
to this basis has a matrix of the form
$\bigoplus_{\lambda\vdash n} f^\lambda W_{\lambda,k}$. The eigenvalues of
$W_{\lambda,k}$ are known to be the contents of the positions occupied by
$k$ in all standard Young tableaux of shape $\lambda$. [\ldots] Darij
Grinberg (private communication) computed the SNF of the matrices
$W_{\lambda,k}$ for $\lambda\vdash n\leq 7$. On the basis of this data we
make the following conjecture.

\textsc{Conjecture 7.4.} \emph{Let $\lambda\vdash n$, $1\leq k\leq n$, and
$W_{\lambda,k}\snf (\alpha_1,\dots,\alpha_{f^\lambda})$. Fix
$1\leq r\leq f^\lambda$. Let $S_r$ be the set of positions $(i,j)$ that $k$
occupies in at least $r$ of the SYT's of shape $\lambda$. Then
$\alpha_{f^\lambda-r+1}=\pm\prod_{(i,j)\in S_r}(j-i)$.}

Note in particular that every SNF diagonal entry is (conjecturally) a
product of some of the eigenvalues of $W_{\lambda,k}$.

For example, when $\lambda=(5,1)$ and $k=5$ we have $f^{(5,1)}=5$ and
$S_1=\{(1,5),(2,1),(1,4)\}$, $S_2=S_3=\{(1,4)\}$, $S_4=S_5=\emptyset$. Hence
$W_{(5,1),5}\snf(1,1,3,3,12)$.
\end{quote}

The omitted passage lists the standard Young tableaux of shape $(5,1)$ and
their eigenvalues. Here $f^\lambda$ is the number of standard Young tableaux
(SYT) of shape $\lambda$, the cell $(i,j)$ is in row $i$ and column $j$, and
$j-i$ is its content. The invariant factors $\alpha_1\mid\alpha_2\mid\cdots$
are taken nonnegative, and an empty product is $1$.

In his comment of December 25, 2024 on~\cite{mo}, Grinberg wrote (the
emphasis, bold in the original, is set in italics here):

\begin{quote}
Note about Question 1: The $W_{\lambda,k}$ in \S7 of Stanley's paper is
defined somewhat clumsily. It should be defined as the matrix representing
the action of the YJM element $X_k$ on the Specht module $\mathcal{S}^\lambda$
over $\mathbb{Z}$ with respect to the standard tableau basis. You can, of
course, take any other basis, since the Smith normal form is
conjugation-invariant. But the direct sum
$\bigoplus_{\lambda \vdash n} f^\lambda W^{\lambda,k}$ is \emph{not} the
action of $X_k$ on the group ring of $S_n$, so you cannot define it through
this sum.
\end{quote}

Question~1 of~\cite{mo}, to which the note refers, reads: ``Is there any
progress on these conjectures ?'' In a second comment Grinberg added that he
was ``still interested in Question 1.''

\paragraph{The reading used.}
We follow Grinberg's clarification. Let $M^\lambda_\Z$ be the free
$\Z$-module on the row tabloids $\{T\}$ of shape $\lambda$, with $\sn$
acting by $\pi\{T\}=\{\pi T\}$, where $\pi T$ replaces each entry $x$ of $T$
by $\pi(x)$. For a tableau $T$ with column stabilizer $C_T$, the polytabloid
is $e_T=\sum_{\pi\in C_T}\operatorname{sgn}(\pi)\,\pi\{T\}$, and the integral
Specht module $S^\lambda_\Z$ is the $\Z$-span of all $e_T$. The standard
polytabloids $e_T$, $T$ standard, form a $\Z$-basis of $S^\lambda_\Z$
(\cite[\S8]{james}, \cite[\S2.6]{sagan}); for the cases we compute, this is
also verified directly by the checkers (Section~\ref{sec:verify}). Then
$W_{\lambda,k}$ is the matrix of $X_k$ in the standard polytabloid basis, and
its Smith form is that of $X_k$ as an endomorphism of $S^\lambda_\Z$, which
does not depend on the basis. By a lattice in $S^\lambda_\Q$ we mean a
full-rank $\Z$-lattice, that is, a $\Z$-submodule of rank $f^\lambda$
spanning $S^\lambda_\Q$; the Smith form of an endomorphism of a lattice is
that of its matrix in any $\Z$-basis. Section~\ref{sec:readings} records
what happens under other readings.

\section{The Jucys--Murphy element \texorpdfstring{$X_3$}{X3}}\label{sec:k3}

\begin{lemma}\label{lem:identity}
In $\Z\sn$, $n\ge3$, let $u=(1,2)$. Then $X_3(X_3-u)=(X_3-u)X_3=2$.
\end{lemma}

\begin{proof}
Write $s=(1,3)$ and $t=(2,3)$, so $X_3=s+t$. Let $c$ and $c^2$ be the two
$3$-cycles on $\{1,2,3\}$. Then $X_3^2=s^2+t^2+st+ts=2+c+c^2$, since $st$ and
$ts=(st)^{-1}$ are the two $3$-cycles. Each of $su$, $tu$, $us$, $ut$ is a
product of two distinct transpositions in $\mathfrak{S}_3$, hence a
$3$-cycle, and $su\ne tu$, $us\ne ut$. So $X_3u=uX_3=c+c^2$, and
$X_3(X_3-u)=(X_3-u)X_3=2$.
\end{proof}

\begin{proof}[Proof of Theorem~\ref{thm:k3}]
Let $L$ be a lattice in $S^\lambda_\Q$ stable under the permutations of
$\{1,2,3\}$. Both $X_3$ and $X_3-u$ lie in $\Z\mathfrak{S}_3$, so in a
$\Z$-basis of $L$ they act by integer matrices $A$ and $B$, and
$AB=BA=2I$ by Lemma~\ref{lem:identity}. Choose unimodular $P,Q$ with
$PAQ=D=\operatorname{diag}(d_1,\dots,d_f)$, $f=f^\lambda$. Then
$D\cdot(Q^{-1}BP^{-1})=2I$ with $Q^{-1}BP^{-1}$ an integer matrix, so
$d_iC_{ii}=2$ for an integer $C_{ii}$, and each $d_i$ is $1$ or $2$. The
number of $d_i$ equal to $2$ is the $2$-adic valuation of $|\det A|$.

The determinant of $X_3$ on $S^\lambda_\Q$ does not depend on $L$. By the
classical theorem of Jucys and Murphy~\cite{jucys,murphy}, $X_3$ acts
diagonally on Young's seminormal basis of $S^\lambda_\Q$, with eigenvalue on
the vector indexed by a standard tableau $T$ equal to the content of the
cell of $3$ in $T$. In a standard tableau the entries $1,2,3$ fill a
partition of $3$, so $3$ occupies $(1,2)$, $(2,1)$, $(1,3)$ or $(3,1)$, of
contents $1,-1,2,-2$. Hence $|\det A|=2^m$, and the Smith form is
$(1^{f-m},2^m)$.
\end{proof}

\begin{proof}[Proof of Corollary~\ref{cor:k3}]
For a cell $c$ let $N_c$ be the number of standard tableaux of shape
$\lambda$ in which $3$ occupies $c$ ($N_c=0$ if $c\notin\lambda$), so that
$m=N_{(1,3)}+N_{(3,1)}$. Every $S_r$ is contained in
$\{(1,2),(2,1),(1,3),(3,1)\}$, and the first two cells have contents $\pm1$.
Hence the prediction is
\[
\alpha_{f-r+1}=2^{[N_{(1,3)}\ge r]+[N_{(3,1)}\ge r]}\qquad(1\le r\le f),
\]
where $[\,\cdot\,]$ is $1$ if the condition holds and $0$ otherwise.

Suppose $\lambda_1\ge3$ and $\lambda$ has at least three rows. The standard
fillings of size three that place $1,2,3$ along the first row, or down the
first column, put $3$ in $(1,3)$, respectively $(3,1)$; both shapes lie in
$\lambda$, and each filling extends to a standard tableau of shape $\lambda$
by filling the remaining cells with $4,\dots,n$ row by row, top to bottom
and left to right. So both contents $2$ and $-2$ occur, $N_{(1,3)}\ge1$ and
$N_{(3,1)}\ge1$, and the prediction gives $\alpha_f=4$. By
Theorem~\ref{thm:k3} every invariant factor divides $2$, so the conjecture
fails, on every lattice of the theorem.

Otherwise at most one of $(1,3)$ and $(3,1)$ lies in $\lambda$, and the
corresponding $N_c$ equals $m$ (if neither does, $m=0$). The prediction is
then $\alpha_{f-r+1}=2$ for $r\le m$ and $1$ for $r>m$, that is,
$(1^{f-m},2^m)$, which is the Smith form given by Theorem~\ref{thm:k3}.
\end{proof}

Theorem~\ref{thm:k3} and Corollary~\ref{cor:k3} are written proofs and are
not machine-checked; the identity of Lemma~\ref{lem:identity} is checked in
Lean. The census of Section~\ref{sec:census} agrees with both for $n\le8$: in
all 63 cases with $k=3$ every invariant factor is $1$ or $2$, there are
exactly $m$ factors equal to $2$, and the disagreements are exactly the
shapes with $\lambda_1\ge3$ and at least three rows.

\begin{remark}
Over $\Z_{(2)}$, the idempotent $P=\frac13(1+c+c^2)$ gives the same
conclusion in another form: $X_3^2=1+3P$, so $X_3$ is invertible on
$(1-P)L_{(2)}$ and acts as $2(1,3)$ on $PL_{(2)}$. Comparing the two
descriptions gives $m=\bigl(f^\lambda+2\chi^\lambda((1,2,3))\bigr)/3$,
where $\chi^\lambda$ is the character of $S^\lambda_\Q$.
\end{remark}

\section{The smallest counterexample}\label{sec:witness}

Let $\lambda=(3,1,1)$, so $n=5$, $f^\lambda=6$ and $X_3=(1,3)+(2,3)$. Order
the standard tableaux as
\[
T_1=\tb{1\,2\,3}{4}{5},\quad
T_2=\tb{1\,2\,4}{3}{5},\quad
T_3=\tb{1\,2\,5}{3}{4},\quad
T_4=\tb{1\,3\,4}{2}{5},\quad
T_5=\tb{1\,3\,5}{2}{4},\quad
T_6=\tb{1\,4\,5}{2}{3},
\]
and let column $j$ of $A$ hold the coordinates of $X_3e_{T_j}$ in the
standard polytabloid basis, so that $X_3e_{T_j}=\sum_i A_{ij}e_{T_i}$. Let
$B=2A^{-1}$; by Lemma~\ref{lem:identity} it is the matrix of $X_3-(1,2)$ in
the same basis, and it is integral. Then
\[
A=\begin{pmatrix}
 2& 0& 0& 0& 0& 0\\
-1&-1& 0& 0& 0& 0\\
 1& 0&-1& 0& 0& 0\\
 0& 1& 0& 1& 0& 0\\
 0& 0& 1& 0& 1& 0\\
 0& 0& 0&-1&-1&-2
\end{pmatrix},\qquad
B=\begin{pmatrix}
 1& 0& 0& 0& 0& 0\\
-1&-2& 0& 0& 0& 0\\
 1& 0&-2& 0& 0& 0\\
 1& 2& 0& 2& 0& 0\\
-1& 0& 2& 0& 2& 0\\
 0&-1&-1&-1&-1&-1
\end{pmatrix},
\]
and $AB=BA=2I$. For example, the
first column of $A$ says $X_3e_{T_1}=2e_{T_1}-e_{T_2}+e_{T_3}$.

\paragraph{The prediction.} The entry $3$ occupies the cell $(1,3)$ in $T_1$,
$(2,1)$ in $T_2$ and $T_3$, $(1,2)$ in $T_4$ and $T_5$, and $(3,1)$ in $T_6$.
Hence $S_1=\{(1,2),(1,3),(2,1),(3,1)\}$ with
$|\prod(j-i)|=|1\cdot2\cdot(-1)\cdot(-2)|=4$, $S_2=\{(1,2),(2,1)\}$ with
product of absolute value $1$, and $S_r=\emptyset$ for $3\le r\le 6$.
Conjecture~7.4 therefore predicts $(1,1,1,1,1,4)$.

\paragraph{The Smith form.} By Theorem~\ref{thm:k3} with $m=2$ ($T_1$ and
$T_6$), the Smith form is $(1,1,1,1,2,2)$. The explicit matrices confirm
this. With
\[
U=\begin{pmatrix}
 0&-1& 0& 0& 0& 0\\
 0&-1&-1& 0& 0& 0\\
 0&-1&-1&-1& 0& 0\\
 0& 1& 1& 1& 1& 0\\
-1& 0& 2& 0& 2& 0\\
 0&-1&-1&-1&-1&-1
\end{pmatrix},\qquad
V=\begin{pmatrix}
 1&-1& 1& 1&-1& 0\\
 0& 1&-1&-1& 1& 0\\
 0& 0& 1& 1&-1& 0\\
 0& 0& 0& 1&-1& 0\\
 0& 0& 0& 0& 1& 0\\
 0& 0& 0& 0& 0& 1
\end{pmatrix},
\]
one has $\det U=-1$, $\det V=1$ and $UAV=\operatorname{diag}(1,1,1,1,2,2)$.
Without transformation matrices: $A$ is lower triangular, so
$\det A=2\cdot(-1)\cdot(-1)\cdot1\cdot1\cdot(-2)=-4$. Modulo $2$, the first row
of $A$ vanishes and the sum of rows $2$ to $6$ vanishes, while rows $2$ to $5$
restricted to columns $2$ to $5$ form a unitriangular matrix; hence $A$ has
rank $4$ over $\F_2$. So exactly two invariant factors are even, their
product with the others is $4$, and the form is $(1,1,1,1,2,2)$. The
predicted form has only one even factor and would require rank $5$ over
$\F_2$.

By Theorem~\ref{thm:k3}, the same Smith form arises on any lattice stable
under the permutations of $\{1,2,3\}$ in a module isomorphic to
$S^{(3,1,1)}_\Q$: for instance the lattices given by polytabloids built from
column tabloids, by a left ideal generated by a Young symmetrizer, by the
dual lattice, or by the saturated lattice $S^\lambda_\Q\cap M^\lambda_\Z$.
No choice of lattice or basis of this kind restores the prediction.

\paragraph{A derivation by hand.}
The matrix can also be obtained without coordinates in the tabloid module.
Let $e_1,\dots,e_5$ be the standard basis of $\Z^5$, with $\mathfrak{S}_5$
permuting it, and let $A_4=\{x\in\Z^5:\sum_i x_i=0\}$. The map $\varphi$
sending the tabloid with $x$ in row $2$ and $y$ in row $3$ to
$e_x\wedge e_y\in\Lambda^2\Z^5$ is $\mathfrak{S}_5$-equivariant. If $T$ has
first column $(c_1,c_2,c_3)$, then $C_T$ is the symmetric group on
$\{c_1,c_2,c_3\}$, and summing the six signed terms gives
\[
\varphi(e_T)=2\,(e_{c_2}-e_{c_1})\wedge(e_{c_3}-e_{c_1})\in 2\Lambda^2A_4 .
\]
For standard $T$ the first column is $(1,a,b)$ with $2\le a<b\le5$, each
pair occurring once, so the six standard polytabloids are sent to $2f_a\wedge
f_b$, where $f_i=e_i-e_1$. These form twice a $\Z$-basis of $\Lambda^2A_4$.
Since $S^\lambda_\Q$ is irreducible and $\varphi\ne0$, the map is injective on
it, and $\frac12\varphi$ is an $\mathfrak{S}_5$-equivariant isomorphism
$S^{(3,1,1)}_\Z\cong\Lambda^2A_4$. Now use the basis $b_i=e_i-e_5$
($1\le i\le4$) of $A_4$ and write $b_{ij}=b_i\wedge b_j$. The transpositions
$(1,3)$ and $(2,3)$ fix $b_4$ and permute $b_1,b_2,b_3$, so $X_3$ preserves
the span of $b_{12},b_{13},b_{23}$ and the span of $b_{14},b_{24},b_{34}$,
with matrices
\[
\begin{pmatrix}0&1&-1\\1&-1&0\\-1&0&-1\end{pmatrix}
\quad\text{and}\quad
\begin{pmatrix}1&0&1\\0&1&1\\1&1&0\end{pmatrix}.
\]
Their determinants are $2$ and $-2$, and each has a $2\times2$ minor equal to
$\pm1$, so each has Smith form $(1,1,2)$. Together they give
$(1,1,1,1,2,2)$ again.

\section{The cases with \texorpdfstring{$n\le5$}{n <= 5}}\label{sec:small}

\begin{proposition}\label{prop:small}
Conjecture~7.4 holds for every $\lambda\vdash n\le4$ and every $1\le k\le n$.
For $n=5$ it fails exactly for $\lambda=(3,1,1)$ with $k=3$ and with $k=5$.
\end{proposition}

For $\lambda=(3,1,1)$ and $k=5$ the Smith form is $(1,1,2,2,4,4)$ and the
prediction $(1,1,1,4,4,4)$. The first statement has a written proof, given
below, which works on every lattice in $S^\lambda_\Q$ preserved by $X_k$.
Both statements are also verified by a standalone checker that recomputes
all 69 cases with $n\le5$ (Section~\ref{sec:verify}), and by the census.

\begin{proof}[Proof for $n\le4$]
We use again that $X_k$ acts diagonally on Young's seminormal basis, with
eigenvalues the contents of the cells occupied by $k$ in the standard
tableaux~\cite{jucys,murphy}. The determinant and the rational rank of
$X_k-a$ ($a\in\Z$) therefore do not depend on the lattice.

If $k=1$, then $X_1=0$, and $1$ occupies $(1,1)$, of content $0$, in every
tableau; the prediction is $(0,\dots,0)$. If $k=2$, then $X_2=(1,2)$ is an
involution, hence unimodular on every lattice, and its Smith form is
$(1,\dots,1)$; the entry $2$ occupies $(1,2)$ or $(2,1)$, of content $\pm1$, so
every predicted factor is $1$. If $f^\lambda=1$, then $\lambda=(n)$ or
$(1^n)$, $X_k$ is the scalar $\pm(k-1)$, and $k$ occupies $(1,k)$ or $(k,1)$,
as predicted. The remaining cases are listed below, with the cells occupied
by $k$ in the standard tableaux (multiplicities as exponents) and the
prediction.
\[
\begin{array}{llll}
\lambda,\ k & \text{cells of }k & \text{eigenvalues} & \text{prediction}\\\hline
(2,1),\ 3 & (1,2),\,(2,1) & 1,-1 & (1,1)\\
(2,2),\ 3 & (1,2),\,(2,1) & 1,-1 & (1,1)\\
(2,2),\ 4 & (2,2)^2 & 0,0 & (0,0)\\
(3,1),\ 3 & (1,3),\,(1,2),\,(2,1) & 2,1,-1 & (1,1,2)\\
(2,1,1),\ 3 & (3,1),\,(1,2),\,(2,1) & -2,1,-1 & (1,1,2)\\
(3,1),\ 4 & (1,3)^2,\,(2,1) & 2,2,-1 & (1,2,2)\\
(2,1,1),\ 4 & (3,1)^2,\,(1,2) & -2,-2,1 & (1,2,2)
\end{array}
\]
In the first two rows $|\det X_k|=1$, so the Smith form is $(1,1)$. In the
third, $X_4$ is diagonalizable with both eigenvalues $0$, so $X_4=0$. In the
fourth and fifth, $|\det X_3|=2$ is squarefree, which forces $(1,1,2)$ (this
also follows from Corollary~\ref{cor:k3}). In the last two, $|\det X_4|=4$,
and $X_4\mp2$ has rational rank $1$, hence rank at most $1$ over $\F_2$; as
$X_4\equiv X_4\mp2\pmod 2$, at least two invariant factors are even, and the
product $4$ forces $(1,2,2)$. Every case agrees with the prediction.
\end{proof}

\section{Other readings}\label{sec:readings}

\paragraph{Other lattices.} For $k=3$, Theorem~\ref{thm:k3} covers every
lattice stable under the permutations of $\{1,2,3\}$. Over the whole census
the following were also compared. The dual lattice
$\operatorname{Hom}(S^\lambda_\Z,\Z)$, with the contragredient action,
carries the transposed matrix, since transpositions are their own inverses;
its Smith form is the same. The saturated lattice
$S^\lambda_\Q\cap M^\lambda_\Z$, which contains $S^\lambda_\Z$, gives the
same Smith form as $S^\lambda_\Z$ in all 135 cases with $n\le6$; it was not
computed for larger $n$. The census also finds that $\lambda$ and its
conjugate $\lambda'$ have the same Smith form for every $k$, in all 416 cases,
as one expects from the classical relation between $S^{\lambda'}$ and the
dual of $S^\lambda$ twisted by the sign character~\cite[\S8]{james}, under
which $X_k$ becomes $-X_k$.

\paragraph{The literal group ring reading.} Read literally, the survey
compares the Smith form of left multiplication by $X_k$ on $\Z\sn$ with that
of a block diagonal matrix holding $f^\lambda$ copies of each predicted list.
Grinberg's comment explains why this is not the intended statement, and it
fails already for $n=4$, $k=4$. Writing $d^m$ for $m$ invariant factors equal
to $d$, left multiplication by $X_4$ on $\Z\mathfrak{S}_4$ has Smith form
$0^4\,1^{16}\,2^4$. The predicted lists are $(3)$ for $(4)$ and $(1^4)$,
$(1,2,2)$ for $(3,1)$ and $(2,1,1)$, and $(0,0)$ for $(2,2)$; with
multiplicities $f^\lambda$ the diagonal holds $0^4\,1^6\,2^{12}\,3^2$, whose
Smith form is $0^4\,1^8\,2^{10}\,6^2$. For $n\le3$ the two agree.

\section{Corrected readings}\label{sec:corrected}

We tested natural modifications of the formula against the same data, the
Smith forms of $X_k$ on $S^\lambda_\Z$ for all $\lambda\vdash n\le8$. Write
$N_c$ for the number of standard tableaux of shape $\lambda$ in which $k$
occupies the cell $c$, so that the formula as stated uses
$S_r=\{c:N_c\ge r\}$. The readings are:
\begin{itemize}\itemsep0pt
\item \emph{as stated}: $\alpha_{f-r+1}=|\prod_{c\in S_r}\mathrm{content}(c)|$;
\item \emph{multiset}: the Smith form of the diagonal matrix of these $f$
products, ignoring their order;
\item \emph{$\alpha_r$ indexing}: the $r$-th product assigned to $\alpha_r$
instead of $\alpha_{f-r+1}$;
\item \emph{exactly $r$, more than $r$, at most $r$}: $S_r$ replaced by
$\{c:N_c=r\}$, $\{c:N_c>r\}$ or $\{c:N_c\le r\}$;
\item \emph{multiplicity classes}: one product for each class of cells with
equal $N_c$, assigned in increasing or decreasing order of $N_c$, or
Smith-normalized;
\item \emph{eigenvalue diagonal}: the Smith form of the diagonal matrix of
all eigenvalues of $X_k$ on $S^\lambda_\Q$;
\item \emph{zero contents dropped}: the formula as stated, with cells of
content $0$ left out of the products;
\item \emph{$X_{k\pm1}$}: the formula for $k$ compared with the Smith form of
$X_{k+1}$ or of $X_{k-1}$;
\item \emph{$Y_k$}: the formula for $k$ compared with the Smith form of
$Y_k=\sum_{i>k}(i,k)$;
\item \emph{$\lambda'$}: the formula for the conjugate shape $\lambda'$
compared with $X_k$ on $S^\lambda$;
\item \emph{fillings of $1,\dots,k$}: $N_c$ replaced by the number of
standard fillings of $1,\dots,k$ inside $\lambda$ with $k$ in $c$;
\item \emph{weak form}: every nonzero invariant factor is a product of some
of the eigenvalues, the remark following the conjecture.
\end{itemize}
A case matches when the predicted list equals the Smith form entry by entry,
taken nonnegative with zeros last. Where a reading has an ordered and a
Smith-normalized version, the table gives the better count. The columns give
the matches among the 240 cases with $n\le7$ and the 416 with $n\le8$ (196
and 350 for $X_{k\pm1}$, where $k\pm1$ must lie in $1,\dots,n$), whether the
reading reproduces Stanley's example $(1,1,3,3,12)$ for $\lambda=(5,1)$,
$k=5$, and the first case, in order of $n$, where it fails.
\[
\begin{array}{lccll}
\text{reading} & n\le7 & n\le8 & (5,1),\,k=5 & \text{first failure}\\\hline
\text{as stated} & 210 & 322 & \text{yes} & (3,1,1),\ k=3\\
\text{multiset} & 210 & 322 & \text{yes} & (3,1,1),\ k=3\\
\lambda' & 210 & 322 & \text{yes} & (3,1,1),\ k=3\\
\alpha_r\text{ indexing} & 122 & 181 & \text{no} & (3,1),\ k=3\\
\text{exactly }r & 96 & 133 & \text{no} & (2,1),\ k=1\\
\text{more than }r & 47 & 70 & \text{no} & (1),\ k=1\\
\text{at most }r & 94 & 131 & \text{no} & (2,1),\ k=1\\
\text{multiplicity classes} & 96 & 133 & \text{no} & (2,1),\ k=1\\
\text{eigenvalue diagonal} & 177 & 279 & \text{yes} & (3,2),\ k=4\\
\text{zero contents dropped} & 173 & 266 & \text{yes} & (3,2),\ k=4\\
X_{k+1} & 3 & 3 & \text{no} & (2),\ k=1\\
X_{k-1} & 3 & 3 & \text{no} & (2),\ k=2\\
Y_k & 26 & 26 & \text{no} & (2),\ k=1\\
\text{fillings of }1,\dots,k & 144 & 203 & \text{yes} & (2,1),\ k=1\\
\text{weak form} & 238 & 410 & \text{yes} & (5,2),\ k=7
\end{array}
\]
The multiset reading agrees with the formula as stated in every case, since
$S_1\supseteq S_2\supseteq\cdots$ makes the stated products a divisibility
chain. The $\lambda'$ reading agrees with it as well: conjugation moves the
counts to transposed cells and changes the sign of the contents, and the
census finds the same Smith form for $\lambda$ and $\lambda'$. So does the
reading with the labels counted from $n$ downward, $Y_{n+1-k}=w_0X_kw_0$ with
$w_0$ the reversal permutation, whose Smith form equals that of $X_k$ in all
416 cases.

No tested reading fits all the data with $n\le7$. The best fit is the formula
as stated, together with its equivalent readings; it matches 210 of the 240
cases with $n\le7$ and first fails at $\lambda=(3,1,1)$, $k=3$. Every other
reading that predicts the whole Smith form first fails at $n\le5$, most of
them at $n\le3$. The weak form, which predicts only a property of it, holds
for all $n\le6$ and fails at $n=7$, for $(5,2)$ and $(2,2,1,1,1)$ with $k=7$
(Section~\ref{sec:census}). We propose no corrected form. The program for
this comparison is included with the optional census programs; it needs the
python-flint library.

\section{The census for \texorpdfstring{$n\le8$}{n <= 8}}\label{sec:census}

We computed the Smith form of $X_k$ on $S^\lambda_\Z$ for every
$\lambda\vdash n\le8$ and every $1\le k\le n$, 416 cases in all, and compared
it with Conjecture~7.4. The numbers of disagreements are:
\[
\begin{array}{l|cccccccc|c}
n & 1&2&3&4&5&6&7&8&\text{total}\\\hline
\text{cases} & 1&4&9&20&35&66&105&176&416\\
\text{disagreements} & 0&0&0&0&2&5&23&64&94
\end{array}
\]
The two disagreements at $n=5$ are those of Proposition~\ref{prop:small}.
Stanley's example $\lambda=(5,1)$, $k=5$, is reproduced: the Smith form is
$(1,1,3,3,12)$, as the survey states. Of the 94 disagreements, 23 have $k=3$
and are accounted for by Corollary~\ref{cor:k3}. Disagreements involve the
prime $2$ in 79 cases and the prime $3$ in 30; they occur both for $k<n$ (85
cases) and for $k=n$ (9 cases), and for two-row shapes from $n=7$ on, the
first being $\lambda=(5,2)$.

The weaker remark following the conjecture, that every SNF entry is a product
of eigenvalues, also fails. For $\lambda=(5,2)$ and $k=7$ the entry $7$
occupies $(1,5)$ or $(2,2)$, so the eigenvalues are $4$ and $0$, yet the Smith
form is $(1,1,1,1,2,4,4,4,4,0,0,0,0,0)$, against the prediction
$(1,1,1,1,1,4,4,4,4,0,0,0,0,0)$; the entry $2$ is not a product of
eigenvalues. The conjugate shape $(2,2,1,1,1)$ behaves in the same way. These
are the first failures of the remark. Nor is the Smith form of $X_k$ always
that of the diagonal matrix of its eigenvalues: this holds in 279 of the 416
cases.

What survives in the data is limited: the conjecture holds for all
$n\le4$ (Proposition~\ref{prop:small}) and at $k=3$ for the shapes of
Corollary~\ref{cor:k3}, and the Smith forms of $\lambda$ and $\lambda'$ agree
in every case. None of the corrected readings of
Section~\ref{sec:corrected} fits the data, and we propose none.

\section{An open point}\label{sec:open}

The documented facts are these. The survey reports that Grinberg computed the
SNF of the matrices $W_{\lambda,k}$ for $\lambda\vdash n\le7$, and that the
conjecture was made on the basis of these data; it defines $W_{\lambda,k}$
through a basis of $\Z\sn$. Grinberg's later comment says that it is defined
``somewhat clumsily'' and names the integral Specht module with its standard basis as the
intended lattice. On that lattice, and by Theorem~\ref{thm:k3} on every
lattice stable under the permutations of $\{1,2,3\}$, the conjecture fails
at $n=5$, and the census finds 30 disagreements with $n\le7$. The data of
the computations reported in the survey are not available to us, and why
they supported the conjecture is not known to us.

\section{Verification}\label{sec:verify}

\paragraph{The interface.} Three kinds of evidence are combined. Lean proves
statements about explicit algebraic objects: the identity of
Lemma~\ref{lem:identity}, and, for the explicit integer matrix $A$, its Smith
form, $AB=2I$, and the unreachability of the predicted form. Standalone
checkers establish that $A$ is the matrix of $X_3$ on $S^{(3,1,1)}_\Z$ in the
standard polytabloid basis, by recomputing it from the definitions. The
general statements, Theorem~\ref{thm:k3}, Corollary~\ref{cor:k3} and the
proof of Proposition~\ref{prop:small} for $n\le4$, are written proofs.

\paragraph{Certificates and checkers.} The witness of
Section~\ref{sec:witness}, and a second witness $\lambda=(5,2)$, $k=7$, are
given as certificates, each holding the standard tableaux in order, the
matrix $A$ of $X_k$, unimodular $U$ and $V$ with $UAV$ diagonal, the Smith
form and the prediction. Two standalone checkers, written independently and
using only the Python standard library, read them; neither trusts the
encoded matrix. The first enumerates the standard
tableaux from all $n!$ fillings, proves that the standard polytabloids form a
$\Z$-basis of the $\Z$-span of all $n!$ polytabloids, recomputes $A$ exactly
from the tabloid module, checks $UAV=\operatorname{diag}(\text{snf})$ with
integer $U,V$ of determinant $\pm1$ and a divisibility chain, and recomputes
the prediction. The second, with separate code, enumerates the standard
tableaux by lattice words, shows that the $\F_2$-span of all polytabloids has
dimension $f^\lambda$, computes the rank of $X_k$ there without reference to
a basis, compares it with the number of odd predicted factors, verifies
$X_ke_{T_j}=\sum_iA_{ij}e_{T_i}$ in the tabloid module for every $j$, and
checks $U$, $V$ and the Smith form with its own arithmetic.
Both reject seven corrupted certificates (a changed entry of $A$; the Smith
form replaced by the prediction; a non-unimodular $U$; a rational $U$ of
determinant $\pm1$; an extra factor; a non-standard basis tableau; an altered
prediction), and both reject the agreeing case $\lambda=(5,1)$, $k=5$, as
showing no discrepancy. A third checker, written separately from both and
importing neither, recomputes every $\lambda\vdash n\le5$ and every $k$: it
proves that the standard polytabloids form a $\Z$-basis, forms the matrix of
$X_k$, computes its Smith form from determinantal divisors, and asserts
Proposition~\ref{prop:small} and Stanley's example $\lambda=(5,1)$, $k=5$.
The entry point is \file{checks/check.py}, with the certificates in
\file{data/}.

\paragraph{Census.} The census was computed by two constructions. The first
generates $S^\lambda_\Z$ inside $M^\lambda_\Z$ as the $\Z\sn$-span of one
polytabloid, closing under simple transpositions with Hermite normal forms;
it uses neither the standard basis nor straightening, and every Smith form
is computed by two implementations. The second uses the standard polytabloid
basis and exact rational linear algebra. They agree in all 135 cases with
$n\le6$. The census output is included as \file{data/census-n8.json}; the
census programs are optional and not part of the check.

\paragraph{Lean.}\sloppy The checked statements are in Lean~4~\cite{lean} with
mathlib~\cite{mathlib}, in the namespace \texttt{StanleySmithForms}:
\texttt{group\_ring\_identity} states the identity of
Lemma~\ref{lem:identity} in the integral group ring of $\sn$, for every
$n\ge3$;
\texttt{A\_det} states $\det A=-4$; \texttt{smith\_form} states
$UAV=\operatorname{diag}(1,1,1,1,2,2)$, and \texttt{U\_det} and
\texttt{V\_det} give $\det U=-1$ and $\det V=1$; \texttt{A\_mul\_B} and
\texttt{B\_mul\_A} state $AB=BA=2I$ for the explicit integer matrix
$B=2A^{-1}$; \texttt{rank\_mod\_two} states that $A$ has rank $4$ over $\F_2$;
\texttt{not\_predicted} states that no integer matrices $P,Q$ with unit
determinant satisfy $PAQ=\operatorname{diag}(1,1,1,1,1,4)$; and
\texttt{counterexample} combines the Smith form and the unreachability of
the predicted form. The passage from the identity, or from $AB=2I$, to the
invariant factors is the written argument of Section~\ref{sec:k3} and is
not formalized. Lean does not check that $A$ is the matrix of $X_3$ on
$S^{(3,1,1)}_\Z$, that is, the construction of the Specht module and its
standard basis, nor that $(1,1,1,1,1,4)$ is the predicted list; the
certificate checkers, which recompute both from $\lambda$ and $k$, and the
derivation of Section~\ref{sec:witness} cover that step.\par\fussy

\begin{thebibliography}{9}

\bibitem{mo}
A.~Chervov, \emph{Conjectures on Smith Normal Form for Jucys-Murphy
elements (Stanley, Grinberg) - status and generalisations?}, with comments
by D.~Grinberg, MathOverflow question 484782, December 25, 2024,
\url{https://mathoverflow.net/questions/484782}.

\bibitem{lean}
L.~de Moura and S.~Ullrich, The Lean~4 theorem prover and programming
language, in \emph{Automated Deduction -- CADE 28}, Lecture Notes in Comput.\
Sci.\ 12699, Springer, 2021, 625--635.

\bibitem{james}
G.~D. James, \emph{The Representation Theory of the Symmetric Groups},
Lecture Notes in Math.\ 682, Springer, Berlin, 1978.

\bibitem{jucys}
A.-A.~A. Jucys, Symmetric polynomials and the center of the symmetric group
ring, \emph{Rep.\ Math.\ Phys.}\ \textbf{5} (1974), 107--112.

\bibitem{mathlib}
The mathlib Community, The Lean mathematical library, in \emph{Proceedings
of the 9th ACM SIGPLAN International Conference on Certified Programs and
Proofs (CPP 2020)}, ACM, 2020, 367--381.

\bibitem{murphy}
G.~E. Murphy, A new construction of Young's seminormal representation of the
symmetric groups, \emph{J. Algebra} \textbf{69} (1981), 287--297.

\bibitem{sagan}
B.~E. Sagan, \emph{The Symmetric Group: Representations, Combinatorial
Algorithms, and Symmetric Functions}, 2nd ed., Grad.\ Texts in Math.\ 203,
Springer, New York, 2001.

\bibitem{stanley}
R.~P. Stanley, Smith normal form in combinatorics, \emph{J. Combin.\ Theory
Ser.~A} \textbf{144} (2016), 476--495; \href{https://arxiv.org/abs/1602.00166}{arXiv:1602.00166}.

\end{thebibliography}
\end{document}
